\documentclass[12pt]{article}
\usepackage[utf8]{inputenc}
\usepackage[T1]{fontenc}
\usepackage{hyperref}
\title{Umpila Language Resources for K-12 Curriculum Development}
\author{Compiled by Claude Sonnet 5 (claude-sonnet-5) for Andr\'as Kornai}
\date{}
\begin{document}
\maketitle
\section{Introduction}
Umpila (ISO 639-3 ump) is a Pama-Nyungan language of Cape York, Queensland, associated with the Lockhart River community. Speaker counts differ. Wikipedia says about 100 people, many elderly. The 2005 National Indigenous Languages Survey, as cited by the State Library of Queensland, reported fewer than a dozen known speakers.
\section{Primary Educational Resources}
A set of readers compiled with the Lockhart River community by Clair Hill and David Thompson is held at the State Library of Queensland. It has six storybooks with audio recordings in both Kuuku Ya'u and Umpila: \url{https://www.slq.qld.gov.au/blog/language-week-week-fifteen-umpila}.
\section{Linguistic and Reference Materials}
The Pama Language Centre says no dictionary is currently available, but it lists several wordlists. These include Hill and Thompson's \emph{Kuuku Ya'u and Umpila Wordlist} (2013), O'Grady's wordlist and Thompson's 1988 outline of Kuuku Ya'u and Umpila: \url{https://www.pamacentre.org.au/umpila/}. Overview: \url{https://en.wikipedia.org/wiki/Umpila_language}.
\section{Religious and Community Resources}
The Pama Language Centre page lists community groups and archival resources, and notes there are no accredited translators.
\section{Digital and Audio Resources}
Audio accompanies the storybook kit above. Another State Library of Queensland post is at \url{https://www.slq.qld.gov.au/blog/2019-international-year-indigenous-languages-word-week-week-fifty-one}.
\section{Curriculum Development Context}
Existing material is a good base: wordlists, readers and audio. There is no dictionary and few speakers, so a curriculum would depend on working with Lockhart River elders and the Pama Language Centre.
\end{document}